\section{Q\&A 与辩题}

\subsection{用户问答速览}

访谈后半段 Burak 回答了听众关心的问题：

\begin{itemize}
    \item \textbf{Q：如何让 agent 在多语言语境中表现一致？}\\
    A：做好 localization pipeline，把 prompt 和 tool manifest 都国际化，使用 multilingual embedding 栈。
    \item \textbf{Q：如何审查 agent 的 hallucination？}\\
    A：结合 search+LLM，生成 answer 时同时输出 source cite，并把 cite 和 user prompt 送入 evaluation harness。
\end{itemize}

\subsection{辩题与反思}

Burak 提出两个值得 debate 的命题：

\begin{enumerate}
    \item Agent 的安全 guardrail 应该先行还是同步迭代？
    \item Multi-modal 能力是否应该优先部署在低风险 task？
\end{enumerate}

他的立场是：guardrail 需要同步迭代（不能等能力再提升才加），multi-modal 应该先在 high-bandwidth observation 不敏感场景验证，再逐渐扩展到高敏感度场景。

\subsection{本章小结}

Q\&A 部分透露了听众最关心的 pain point，也凸显出企业要在能力与 governance 之间持续 debate。

